%% kapa: exam=1 round=33 year=2022 period=1 form=A subject=economics subject_ko=경제학원론 seq=29 number=69
%% source: Q-Net 감정평가사 기출문제 (cst00309) artlSeq=5209997 file=raw/1008/5209997/2022년도 제33회 감정평가사 1차시험 1교시(민법,경제학,부동산학).pdf
%% conversion: kordoc 4.24.51 -> tools/gen1.py (auto); review=proofread-sonnet; flags=none
\begin{question}{69}
폐쇄경제 IS-LM 모형에서 재정정책과 통화정책이 생산량에 미치는 효과의 크기에 관한 설명으로 옳은 것을 모두 고른 것은? (단, IS는 우하향, LM은 우상향하는 직선이다.)
\begin{framedtext}
ㄱ. 투자가 이자율에 민감할수록 통화정책의 효과가 작다. \par
ㄴ. 화폐수요가 이자율에 민감할수록 재정정책의 효과가 크다. \par
ㄷ. 한계소비성향이 클수록 통화정책의 효과가 크다.
\end{framedtext}
\begin{choices}
  \item ㄱ
  \item ㄷ
  \item ㄱ, ㄴ
  \item ㄱ, ㄷ
  \item ㄴ, ㄷ
\end{choices}
\end{question}
\answer{5}
